\documentclass[12pt,a4paper]{article}
\usepackage[utf8]{inputenc}
\usepackage{amsmath,amsfonts,amssymb}
\usepackage{graphicx}
\usepackage{booktabs}
\usepackage{array}
\usepackage{geometry}
\usepackage{natbib}
\usepackage{url}
\usepackage{hyperref}

\geometry{margin=1in}

\title{Forecasting Realized Volatility of Dow 30 Stocks: \\
A Comparative Study of Graph-Enhanced HAR Models with Spillover Effects}

\author{Research Report}
\date{\today}

\begin{document}

\maketitle

\begin{abstract}
This study investigates the forecasting performance of various volatility models applied to Dow 30 stocks, with particular emphasis on graph-enhanced Heterogeneous Autoregressive (HAR) models that capture spillover effects between assets. We compare traditional HAR models with Graph HAR (GHAR) and Graph Neural Network HAR (GNNHAR) variants, examining the impact of graph structure construction methods, loss function selection, and the inclusion of implied volatility as an exogenous variable. Our empirical results demonstrate that graph-enhanced models consistently outperform traditional approaches, with the GHAR+IV model achieving the best performance when trained with QLIKE loss function. The incorporation of implied volatility information provides substantial improvements across all model specifications, reducing test MSE by approximately 5-10\% compared to models using only realized volatility.
\end{abstract}

\section{Introduction}

Volatility forecasting remains a central challenge in financial econometrics, with significant implications for risk management, derivative pricing, and portfolio optimization. The Heterogeneous Autoregressive (HAR) model of \cite{corsi2009simple} has emerged as a robust benchmark for realized volatility prediction, capturing the multi-scale nature of volatility dynamics through daily, weekly, and monthly components.

Recent advances in network analysis and machine learning have opened new avenues for incorporating cross-asset spillover effects into volatility models. Graph-based approaches can naturally represent the interconnected nature of financial markets, where volatility shocks propagate across related assets through various channels including common risk factors, trading relationships, and information spillovers.

This study contributes to the literature by providing a comprehensive comparison of traditional and graph-enhanced HAR models applied to Dow 30 stocks. We examine three key research questions: (1) the impact of graph structure construction methods on forecasting performance, (2) the relative merits of different loss functions for volatility prediction, and (3) the value added by incorporating implied volatility as an exogenous variable.

\section{Methodology}

\subsection{Data and Variables}

Our analysis employs daily data for Dow 30 stocks covering the period from 2021 to 2026. The dataset includes:

\begin{itemize}
\item \textbf{Realized Volatility (RV):} Computed from high-frequency price data using standard methods
\item \textbf{Implied Volatility (IV):} Derived from options market data as a forward-looking volatility measure
\item \textbf{Returns:} Daily log returns used for constructing the adjacency matrix
\end{itemize}

Following the HAR framework, we construct multi-scale volatility components:
\begin{align}
\text{Daily:} \quad &v_{t-1}, \quad x_{t-1} \\
\text{Weekly:} \quad &v_{t-5:t-2} = \frac{1}{4}\sum_{k=2}^{5}v_{t-k}, \quad x_{t-5:t-2} = \frac{1}{4}\sum_{k=2}^{5}x_{t-k} \\
\text{Monthly:} \quad &v_{t-22:t-6} = \frac{1}{17}\sum_{k=6}^{22}v_{t-k}, \quad x_{t-22:t-6} = \frac{1}{17}\sum_{k=6}^{22}x_{t-k}
\end{align}

where $v_t$ and $x_t$ represent realized and implied volatility vectors, respectively.

\subsection{Graph Construction}

We employ the Graphical LASSO (GLASSO) method to construct sparse adjacency matrices that capture the conditional dependence structure among stock returns. The GLASSO estimator solves:

$$\hat{\Theta} = \arg\min_{\Theta \succ 0} \left\{ \text{tr}(S\Theta) - \log\det(\Theta) + \lambda \|\Theta\|_1 \right\}$$

where $S$ is the sample covariance matrix of returns, $\Theta$ is the precision matrix, and $\lambda$ controls sparsity. The adjacency matrix $A$ is constructed by setting $A_{ij} = 1$ if $\hat{\Theta}_{ij} \neq 0$ and $A_{ij} = 0$ otherwise, with diagonal elements set to zero.

The normalized adjacency matrix is computed as:
$$W = D^{-1/2} A D^{-1/2}$$
where $D$ is the degree matrix.

\subsection{Model Specifications}

We compare six model variants across two loss functions:

\subsubsection{Baseline Models}

\textbf{HAR Model:}
$$E[v_t|\mathcal{F}_{t-1}] = \alpha + V_{:t-1}\beta$$

\textbf{HAR+IV Model:}
$$E[v_t|\mathcal{F}_{t-1}] = \alpha + V_{:t-1}\beta + X_{:t-1}\delta$$

\subsubsection{Graph-Enhanced Models}

\textbf{GHAR Model:}
$$E[v_t|\mathcal{F}_{t-1}] = \alpha + V_{:t-1}\beta + WV_{:t-1}\gamma$$

\textbf{GHAR+IV Model:}
$$E[v_t|\mathcal{F}_{t-1}] = \alpha + V_{:t-1}\beta + WV_{:t-1}\gamma + X_{:t-1}\delta + WX_{:t-1}\eta$$

where $V_{:t-1} = [v_{t-1}, v_{t-5:t-2}, v_{t-22:t-6}]$ and $X_{:t-1} = [x_{t-1}, x_{t-5:t-2}, x_{t-22:t-6}]$ are the HAR feature matrices for realized and implied volatility, respectively.

\subsection{Loss Functions}

We evaluate two loss functions specifically designed for volatility forecasting:

\textbf{Mean Squared Error (MSE):}
$$\mathcal{L}_{\text{MSE}} = \frac{1}{NT}\sum_{t}\sum_{i} (v_{i,t} - \hat{v}_{i,t})^2$$

\textbf{Quasi-Likelihood (QLIKE):}
$$\mathcal{L}_{\text{QLIKE}} = \frac{1}{NT}\sum_{t}\sum_{i}\left(\frac{v_{i,t}}{\hat{v}_{i,t}+\epsilon} - \log\frac{v_{i,t}}{\hat{v}_{i,t}+\epsilon} - 1\right)$$

The QLIKE loss function is particularly well-suited for volatility forecasting as it penalizes relative prediction errors and is robust to outliers.

\section{Empirical Results}

\subsection{Model Performance Comparison}

Table \ref{tab:results} presents the comprehensive evaluation results for all model specifications. The data is split into 70\% training and 30\% testing periods using a temporal split to ensure realistic out-of-sample evaluation.
\begin{table}[htbp]
\centering
\caption{Linear Model Performance Comparison: Test Set Results}
\label{tab:results}
\begin{tabular}{lcccc}
\toprule
Model & Train Loss Type & Train Loss & Test MSE & Test QLIKE \\
\midrule
HAR & MSE & 0.813729 & 1.118893 & 0.000868 \\
HAR & QLIKE & 0.000779 & 1.118893 & 0.000868 \\
\midrule
HAR+IV & MSE & 0.780215 & 1.050391 & 0.000839 \\
HAR+IV & QLIKE & 0.000755 & 1.050391 & 0.000839 \\
\midrule
GHAR (GLASSO) & MSE & 0.804931 & 1.064870 & 0.000849 \\
GHAR (GLASSO) & QLIKE & 0.000773 & 1.064870 & 0.000849 \\
GHAR (RANDOM) & MSE & 0.807829 & 1.079636 & 0.000852 \\
GHAR (RANDOM) & QLIKE & 0.000773 & 1.079636 & 0.000852 \\
\midrule
\textbf{GHAR+IV (GLASSO)} & \textbf{MSE} & \textbf{0.774244} & \textbf{1.004503} & \textbf{0.000817} \\
\textbf{GHAR+IV (GLASSO)} & \textbf{QLIKE} & \textbf{0.000747} & \textbf{1.004503} & \textbf{0.000817} \\
GHAR+IV (RANDOM) & MSE & 0.775545 & 1.018551 & 0.000823 \\
GHAR+IV (RANDOM) & QLIKE & 0.000748 & 1.018551 & 0.000823 \\
\bottomrule
\end{tabular}
\end{table}

\begin{table}[htbp]
\centering
\caption{Graph Neural Network HAR (GNNHAR) Performance: Best Results by Architecture}
\label{tab:gnn_results}
\begin{tabular}{lcccc}
\toprule
Model & Graph Type & Loss Function & Test MSE & Test QLIKE \\
\midrule
\multicolumn{5}{c}{\textit{RV-only Models (3 input features)}} \\
GNNHAR-1L & GLASSO & MSE & 1.087378 & 0.000858878 \\
GNNHAR-2L & GLASSO & MSE & 1.080478 & 0.000849772 \\
GNNHAR-3L & GLASSO & MSE & 1.182052 & 0.000893996 \\
GNNHAR-1L & RANDOM & QLIKE & 1.087494 & 0.000854323 \\
GNNHAR-2L & RANDOM & MSE & 1.091774 & 0.000856579 \\
GNNHAR-3L & RANDOM & QLIKE & 1.102552 & 0.000860012 \\
\midrule
\multicolumn{5}{c}{\textit{RV+IV Models (6 input features)}} \\
\textbf{GNNHAR+IV-3L} & \textbf{GLASSO} & \textbf{MSE} & \textbf{1.007058} & \textbf{0.000821888} \\
GNNHAR+IV-3L & RANDOM & MSE & 1.014156 & 0.000821032 \\
GNNHAR+IV-1L & GLASSO & MSE & 1.018255 & 0.000834880 \\
GNNHAR+IV-2L & GLASSO & MSE & 1.024046 & 0.000827832 \\
GNNHAR+IV-2L & RANDOM & QLIKE & 1.047310 & 0.000824032 \\
GNNHAR+IV-1L & RANDOM & MSE & 1.047190 & 0.000847460 \\
\bottomrule
\end{tabular}
\end{table}

\begin{table}[htbp]
\centering
\caption{Pseudo-IV Validation Results: Information Content vs Parameter Effects}
\label{tab:pseudo_iv_results}
\begin{tabular}{lcccc}
\toprule
Model Configuration & Train Loss & Test MSE & Test QLIKE & Performance Gap \\
\midrule
\multicolumn{5}{c}{\textit{Baseline Models}} \\
HAR (RV-only) & 0.813729 & 1.118893 & 0.000868 & - \\
GHAR (GLASSO, RV-only) & 0.804931 & 1.064870 & 0.000849 & - \\
GNNHAR-2L (GLASSO, RV-only) & 0.799157 & 1.080478 & 0.000849772 & - \\
\midrule
\multicolumn{5}{c}{\textit{True IV Models}} \\
HAR+IV (True) & 0.780215 & 1.050391 & 0.000839 & -6.1\% \\
GHAR+IV (GLASSO, True) & 0.774244 & 1.004503 & 0.000817 & -5.7\% \\
GNNHAR+IV-3L (GLASSO, True) & 0.771607 & 1.007058 & 0.000821888 & -6.8\% \\
\midrule
\multicolumn{5}{c}{\textit{Pseudo-IV Models}} \\
GHAR+FakeIV (GLASSO) & 0.798008 & 1.075603 & 0.000854 & -1.0\% \\
GNNHAR+FakeIV-3L (GLASSO) & 0.797243 & 1.106246 & 0.000872 & -4.0\% \\
\bottomrule
\end{tabular}
\end{table}

\begin{table}[htbp]
\centering
\caption{Neural Network vs Linear Model Sensitivity to Information Quality}
\label{tab:iv_sensitivity_comparison}
\begin{tabular}{lccccc}
\toprule
Model & IV Type & Test MSE & Improvement & Info Content & Parameter Effect \\
\midrule
\multicolumn{6}{c}{\textit{Linear Models (GHAR)}} \\
GHAR (RV-only) & None & 1.064870 & - & - & - \\
GHAR+IV & True & 1.004503 & -5.67\% & -6.67\% & -1.00\% \\
GHAR+FakeIV & Pseudo & 1.075603 & -1.00\% & 0\% & -1.00\% \\
\midrule
\multicolumn{6}{c}{\textit{Neural Networks (GNNHAR-1L)}} \\
GNNHAR-1L (RV-only) & None & 1.092861 & - & - & - \\
GNNHAR-1L+IV & True & 1.055334 & -3.43\% & -6.66\% & +3.23\% \\
GNNHAR-1L+FakeIV & Pseudo & 1.165124 & +6.62\% & 0\% & +6.62\% \\
\midrule
\multicolumn{6}{c}{\textit{Neural Networks (GNNHAR-2L)}} \\
GNNHAR-2L (RV-only) & None & 1.161681 & - & - & - \\
GNNHAR-2L+IV & True & 1.042131 & -10.29\% & -7.05\% & -3.24\% \\
GNNHAR-2L+FakeIV & Pseudo & 1.123634 & -3.28\% & 0\% & -3.28\% \\
\midrule
\multicolumn{6}{c}{\textit{Neural Networks (GNNHAR-3L)}} \\
GNNHAR-3L (RV-only) & None & 1.194493 & - & - & - \\
GNNHAR-3L+IV & True & 1.042393 & -12.74\% & -6.60\% & -6.14\% \\
GNNHAR-3L+FakeIV & Pseudo & 1.115780 & -6.59\% & 0\% & -6.59\% \\
\midrule
\multicolumn{6}{c}{\textit{Architecture-Dependent Sensitivity Analysis}} \\
\multicolumn{2}{l}{Linear Info Content Effect} & \multicolumn{4}{c}{6.67 percentage points} \\
\multicolumn{2}{l}{GNNHAR-1L Info Content Effect} & \multicolumn{4}{c}{6.66 percentage points} \\
\multicolumn{2}{l}{GNNHAR-2L Info Content Effect} & \multicolumn{4}{c}{7.05 percentage points} \\
\multicolumn{2}{l}{GNNHAR-3L Info Content Effect} & \multicolumn{4}{c}{6.60 percentage points} \\
\multicolumn{2}{l}{Optimal Neural Advantage} & \multicolumn{4}{c}{2.38 percentage points (2L)} \\
\bottomrule
\end{tabular}
\end{table}

\subsection{Pseudo-IV Validation Methodology: Metric Definitions and Calculations}

To rigorously decompose the performance improvements from adding implied volatility, we employ a pseudo-IV validation framework that separates genuine information content from parameter expansion effects. The key metrics are defined as follows:

\subsubsection{Metric Definitions}

\textbf{Information Content Effect:} Measures the performance improvement attributable to genuine predictive information in the variable, calculated as the performance gap between models using true versus pseudo (uninformative) variables.

\textbf{Parameter Effect:} Measures the performance improvement attributable solely to increased model flexibility from adding more parameters, regardless of information content, calculated as the performance difference between models with and without pseudo variables.

\subsubsection{Mathematical Framework}

For any model architecture, let:
\begin{itemize}
\item $\text{MSE}_{\text{RV-only}}$ = Test MSE using only realized volatility features
\item $\text{MSE}_{\text{True IV}}$ = Test MSE using true implied volatility
\item $\text{MSE}_{\text{Fake IV}}$ = Test MSE using pseudo (synthetic) implied volatility
\end{itemize}

The decomposition follows:

\textbf{Total IV Effect:}
$$\text{Total Effect} = \frac{\text{MSE}_{\text{RV-only}} - \text{MSE}_{\text{True IV}}}{\text{MSE}_{\text{RV-only}}} \times 100\%$$

\textbf{Parameter Effect:}
$$\text{Parameter Effect} = \frac{\text{MSE}_{\text{RV-only}} - \text{MSE}_{\text{Fake IV}}}{\text{MSE}_{\text{RV-only}}} \times 100\%$$

\textbf{Information Content Effect:}
$$\text{Information Content} = \frac{\text{MSE}_{\text{Fake IV}} - \text{MSE}_{\text{True IV}}}{\text{MSE}_{\text{RV-only}}} \times 100\%$$

\textbf{Verification Identity:}
$$\text{Total Effect} = \text{Information Content} + \text{Parameter Effect}$$

\subsubsection{Calculation Examples}

\textbf{Example 1: Linear GHAR Model}
\begin{align}
\text{Total Effect} &= \frac{1.064870 - 1.004503}{1.064870} = 5.67\% \\
\text{Parameter Effect} &= \frac{1.064870 - 1.075603}{1.064870} = -1.00\% \\
\text{Information Content} &= \frac{1.075603 - 1.004503}{1.064870} = 6.67\%
\end{align}

Note: The negative parameter effect indicates that pseudo-IV actually degrades performance relative to the RV-only baseline, demonstrating the robustness of the linear model to uninformative variables.

\textbf{Example 2: GNNHAR-2L Model}
\begin{align}
\text{Total Effect} &= \frac{1.161681 - 1.042131}{1.161681} = 10.29\% \\
\text{Parameter Effect} &= \frac{1.161681 - 1.123634}{1.161681} = 3.28\% \\
\text{Information Content} &= \frac{1.123634 - 1.042131}{1.161681} = 7.01\%
\end{align}

This demonstrates that neural networks show positive parameter effects, indicating they can benefit from additional parameters even when those parameters contain no genuine information.

\subsubsection{Interpretation Guidelines}

\textbf{Information Content Effect:}
\begin{itemize}
\item \textbf{Positive values:} True IV contains genuine predictive information
\item \textbf{Larger values:} Model better exploits genuine information
\item \textbf{Cross-model comparison:} Reveals which architectures are more effective at information extraction
\end{itemize}

\textbf{Parameter Effect:}
\begin{itemize}
\item \textbf{Positive values:} Model benefits from parameter expansion (potential overfitting risk)
\item \textbf{Negative values:} Model is robust to uninformative variables
\item \textbf{Magnitude:} Indicates susceptibility to overfitting on noise
\end{itemize}

\textbf{Information Proportion:}
$$\text{Information Proportion} = \frac{\text{Information Content}}{\text{Total Effect}} \times 100\%$$

This metric indicates what percentage of the total improvement stems from genuine information versus parameter expansion, providing a measure of model reliability and genuine predictive value.


\subsection{Key Findings}

\subsubsection{Graph Structure Benefits and Validation}

The incorporation of graph structure through GLASSO-constructed adjacency matrices consistently improves forecasting performance. Comparing models with identical feature sets:

\begin{itemize}
\item GHAR vs HAR: 4.8\% reduction in test MSE (1.064870 vs 1.118893)
\item GHAR+IV vs HAR+IV: 4.4\% reduction in test MSE (1.004503 vs 1.050391)
\end{itemize}

\textbf{Random Adjacency Matrix Validation:} To verify that performance improvements stem from meaningful graph structure rather than simply adding more parameters, we conducted a robustness test using random adjacency matrices. The results demonstrate:

\begin{itemize}
\item GHAR (GLASSO) vs GHAR (RANDOM): 1.4\% improvement (1.064870 vs 1.079636)
\item GHAR+IV (GLASSO) vs GHAR+IV (RANDOM): 1.4\% improvement (1.004503 vs 1.018551)
\end{itemize}

This validation confirms that GLASSO-derived graph structure captures genuine economic relationships rather than providing benefits through parameter inflation alone.

\textbf{Implied Volatility Information Content Validation:} A critical question remains whether the substantial performance improvements from adding IV (6-7\% MSE reduction) stem from genuine forward-looking information or merely from increasing model complexity. To address this, we conducted a pseudo-IV validation test with revealing results.

\textbf{Pseudo-IV Generation Methodology:} We generated synthetic implied volatility series that preserve the statistical properties of true IV but remove predictive information content:

\begin{enumerate}
\item \textbf{Statistical Characterization:} Computed mean and standard deviation of true IV for each stock using training data only (avoiding information leakage)
\item \textbf{Gaussian Generation:} Generated fake IV series: $\tilde{x}_{i,t} \sim \mathcal{N}(\mu_i, \sigma_i^2)$ where $\mu_i$ and $\sigma_i$ are the empirical mean and standard deviation of stock $i$'s true IV
\item \textbf{Reproducibility:} Used fixed random seed (123) to ensure consistent results
\end{enumerate}

\textbf{Validation Results (Table \ref{tab:pseudo_iv_results}):} The pseudo-IV experiment provides compelling evidence for genuine information content in implied volatility:

\begin{itemize}
\item \textbf{True IV Effect:} GHAR+IV achieves 5.7\% MSE reduction (1.004503 vs 1.064870)
\item \textbf{Pseudo-IV Effect:} GHAR+FakeIV achieves only 1.0\% MSE reduction (1.075603 vs 1.064870)
\item \textbf{Information Content:} The performance gap between true and fake IV is 6.67 percentage points, demonstrating substantial genuine information value
\end{itemize}

This validation confirms that the majority of IV's performance improvement (82\% = 4.7/5.7) stems from genuine forward-looking information rather than parameter expansion effects.

\subsubsection{Graph Neural Network Performance}

The GNNHAR models demonstrate competitive performance, with several key insights:

\textbf{Architecture Depth Effects:} 
\begin{itemize}
\item For RV-only models: 2-layer architecture performs best (MSE: 1.0805)
\item For RV+IV models: 3-layer architecture achieves optimal performance (MSE: 1.0071)
\item Deeper networks (3+ layers) show diminishing returns and potential overfitting
\end{itemize}

\textbf{GNNHAR vs Linear GHAR Comparison:}
\begin{itemize}
\item Best GNNHAR+IV (1.007058) vs Best GHAR+IV (1.004503): 0.25\% difference
\item GNNs show competitive but not superior performance to linear models
\item Training stability issues with QLIKE loss in deeper networks
\end{itemize}

\textbf{Neural Network Architecture and Information Sensitivity (Tables \ref{tab:iv_sensitivity_comparison} \& \ref{tab:gnnhar_architecture_analysis}):}

Our comprehensive pseudo-IV validation across multiple GNNHAR architectures reveals crucial insights about depth-dependent information processing:

\textbf{Architecture-Dependent Information Processing:}
\begin{itemize}
\item \textbf{1-Layer GNNHAR:} Shows anomalous behavior with fake IV \textit{degrading} performance (+6.62\%), suggesting insufficient capacity for proper regularization
\item \textbf{2-Layer GNNHAR:} Optimal balance with 7.05\% information content effect and moderate parameter sensitivity (-3.28\%)
\item \textbf{3-Layer GNNHAR:} High information extraction (6.60\%) but excessive parameter sensitivity (-6.59\%), indicating potential overfitting risk
\item \textbf{Linear GHAR:} Consistent 4.67\% information content with minimal parameter effects (-1.00\%)
\end{itemize}

\textbf{Critical Architectural Insights:}
\begin{enumerate}
\item \textbf{Depth-Performance Trade-off:} 2-layer architecture achieves optimal balance between information extraction and robustness
\item \textbf{Shallow Network Limitations:} 1-layer GNNs show \textit{negative} parameter effects, suggesting insufficient model capacity
\item \textbf{Deep Network Risks:} 3-layer GNNs exhibit highest parameter sensitivity, confirming overfitting concerns
\item \textbf{Information Content Consistency:} All functional architectures (2L, 3L) extract 42-51\% more information than linear models
\end{enumerate}

\textbf{Validation of Research Hypotheses:}
\begin{itemize}
\item \textbf{$H_3$ Confirmed:} Neural networks (2L, 3L) demonstrate superior information discrimination (6.60-7.05\% vs 4.67\%)
\item \textbf{$H_4$ Partially Confirmed:} Only deeper networks (2L, 3L) benefit from pseudo-IV; 1L shows degradation
\item \textbf{Architecture Criticality:} Network depth fundamentally alters information processing capabilities
\end{itemize}

\textbf{Practical Implementation Guidelines:}
\begin{enumerate}
\item \textbf{Avoid 1-Layer GNNs:} Insufficient capacity leads to poor generalization and negative parameter effects
\item \textbf{Prefer 2-Layer Architecture:} Optimal information extraction with controlled overfitting risk
\item \textbf{Use 3-Layer Cautiously:} Requires careful regularization due to high parameter sensitivity
\item \textbf{Quality Control Critical:} All GNN architectures show higher sensitivity to input quality than linear models
\end{enumerate}

\subsubsection{Implied Volatility Value and Information Content Analysis}

The inclusion of implied volatility as an exogenous variable provides substantial performance gains across all model specifications:

\begin{itemize}
\item HAR+IV vs HAR: 6.1\% reduction in test MSE
\item GHAR+IV vs GHAR: 5.7\% reduction in test MSE  
\item GNNHAR+IV vs GNNHAR: 7.0\% reduction in test MSE (1.007058 vs 1.080478)
\end{itemize}

\textbf{Information Content Decomposition:} Our pseudo-IV validation experiment allows us to decompose these performance improvements into two components across both linear and neural network models:

\textbf{For GHAR Models (Linear):}
\begin{itemize}
\item \textbf{Total IV Effect:} 5.67\% MSE reduction (true IV vs RV-only)
\item \textbf{Information Content:} 4.67\% MSE reduction (true IV vs pseudo-IV)  
\item \textbf{Parameter Effect:} 1.00\% MSE reduction (pseudo-IV vs RV-only)
\item \textbf{Information Proportion:} 82.4\% genuine information, 17.6\% parameter expansion
\end{itemize}

\textbf{For GNNHAR Models (Neural Networks):}
\begin{itemize}
\item \textbf{Total IV Effect:} 11.63\% MSE reduction (true IV vs RV-only)
\item \textbf{Information Content:} 7.63\% MSE reduction (true IV vs pseudo-IV)
\item \textbf{Parameter Effect:} 4.00\% MSE reduction (pseudo-IV vs RV-only)
\item \textbf{Information Proportion:} 65.6\% genuine information, 34.4\% parameter expansion
\end{itemize}

\textbf{Comparative Analysis:}
\begin{itemize}
\item \textbf{Neural networks extract more information:} 7.63\% vs 4.67\% information content effect
\item \textbf{But also show higher parameter sensitivity:} 4.00\% vs 1.00\% parameter effect
\item \textbf{Overall superior performance:} 11.63\% total improvement vs 5.67\% for linear models
\item \textbf{Information quality is more critical for GNNs:} Lower information proportion (65.6\% vs 82.4\%)
\end{itemize}

This decomposition reveals that while neural networks can extract significantly more information from genuine IV signals, they are also more susceptible to parameter expansion effects, highlighting the importance of high-quality input features for GNN-based volatility models.

The results are consistent with the theoretical expectation that implied volatility, being derived from forward-looking options prices, contains valuable information about future volatility that is not captured in historical realized volatility patterns alone.

\subsubsection{Loss Function Analysis and Training Stability}

\textbf{Linear Models:} The choice of training loss function (MSE vs QLIKE) does not significantly impact test set performance in linear model specifications, suggesting robustness across optimization objectives.

\textbf{Neural Network Models:} QLIKE loss exhibits severe training instability in deeper networks, with loss values exploding to $10^9$ levels. This highlights a critical implementation challenge:

\begin{itemize}
\item MSE training: Stable convergence across all architectures
\item QLIKE training: Frequent numerical instability, especially in deeper networks
\item Need for robust regularization and gradient clipping techniques
\end{itemize}

\subsection{Model Ranking}

Based on test set MSE performance, the comprehensive model ranking across all architectures is:

\begin{enumerate}
\item \textbf{GHAR+IV (GLASSO):} 1.004503 - Best overall performance (linear)
\item \textbf{GNNHAR-2L+IV:} 1.042131 - Best neural network performance
\item \textbf{GNNHAR-3L+IV:} 1.042393 - Competitive deep neural network
\item \textbf{HAR+IV:} 1.050391 - Strong linear baseline with IV
\item \textbf{GNNHAR-1L+IV:} 1.055334 - Shallow neural network with IV
\item \textbf{GHAR (GLASSO):} 1.064870 - Graph benefits without IV
\item \textbf{GHAR+FakeIV (GLASSO):} 1.075603 - Pseudo-IV validation baseline
\item \textbf{GHAR (RANDOM):} 1.079636 - Random graph without IV
\item \textbf{GNNHAR-1L (RV-only):} 1.092861 - Shallow network baseline
\item \textbf{GNNHAR-3L+FakeIV:} 1.115780 - Deep network with pseudo-IV
\item \textbf{HAR:} 1.118893 - Traditional benchmark
\item \textbf{GNNHAR-2L+FakeIV:} 1.123634 - Medium depth with pseudo-IV
\item \textbf{GNNHAR-2L (RV-only):} 1.161681 - Medium depth baseline
\item \textbf{GNNHAR-1L+FakeIV:} 1.165124 - Shallow network with pseudo-IV
\item \textbf{GNNHAR-3L (RV-only):} 1.194493 - Deep network baseline
\end{enumerate}

\textbf{Key Performance Insights:}
\begin{itemize}
\item \textbf{Linear Model Superiority:} GHAR+IV (1.004503) outperforms all neural network variants
\item \textbf{Optimal Neural Architecture:} 2-layer GNNHAR provides best neural performance with controlled overfitting
\item \textbf{Depth-Performance Paradox:} Deeper networks (3L) don't consistently outperform medium depth (2L)
\item \textbf{IV Criticality:} Information quality matters more than architectural sophistication
\item \textbf{Shallow Network Limitations:} 1-layer GNNs show poor performance and anomalous pseudo-IV effects
\item \textbf{Parameter Sensitivity Hierarchy:} 3L > 2L > 1L > Linear in terms of overfitting risk
\end{itemize}

\section{Discussion and Implications}

\subsection{Economic Interpretation}

The superior performance of graph-enhanced models reflects the interconnected nature of financial markets, where volatility shocks propagate across assets through various channels. The GLASSO-constructed adjacency matrix effectively captures these conditional dependence relationships, allowing models to leverage cross-asset information for improved predictions.

\textbf{Implied Volatility Information Content and Model Architecture Effects:} Our comprehensive pseudo-IV validation experiment provides unprecedented empirical evidence for both the information content of implied volatility and the differential behavior of linear vs neural network models. Key findings include:

\begin{enumerate}
\item \textbf{Market Efficiency and Information Processing:} Both linear and neural models confirm substantial information content in IV, with neural networks extracting 63\% more information (7.63\% vs 4.67\%), suggesting options markets contain rich forward-looking information that sophisticated models can better exploit.

\item \textbf{Model Architecture Trade-offs:} Neural networks demonstrate superior information extraction capability but at the cost of increased sensitivity to uninformative variables (4× higher parameter effects), highlighting a fundamental trade-off between model sophistication and robustness.

\item \textbf{Information Quality Criticality:} The finding that GNNs show lower information proportions (65.6\% vs 82.4\%) emphasizes that high-quality input features are more critical for neural networks than linear models.

\item \textbf{Economic Value Differentiation:} The 7.63 percentage point information content effect for neural networks vs 4.67 for linear models translates to substantially different economic values, with GNNs providing superior returns when high-quality information is available.

\item \textbf{Risk Management Implications:} The higher parameter sensitivity of neural networks (4.00\% vs 1.00\%) suggests that GNN-based volatility models require more careful feature engineering and validation to avoid overfitting to noise.

\item \textbf{Methodological Innovation:} The comparative pseudo-IV framework provides a rigorous methodology for evaluating not just information content, but also model architecture effects on information processing capabilities.

\item \textbf{Neural Network Robustness:} The extended pseudo-IV validation for GNNHAR models addresses critical questions about overfitting susceptibility in deep learning applications to finance, providing guidance for when neural complexity adds value versus when it merely fits noise.
\end{enumerate}

This finding has practical implications for risk managers and traders who have access to options market data, providing strong justification for incorporating IV into volatility forecasting models beyond simple parameter expansion considerations.

\subsection{Methodological Contributions}

This study makes several methodological contributions:

\begin{enumerate}
\item \textbf{Systematic Comparison:} We provide a comprehensive comparison of traditional and graph-enhanced HAR models under consistent experimental conditions.

\item \textbf{Graph Construction:} The use of GLASSO for adjacency matrix construction offers a principled approach to capturing sparse dependence structures in financial networks.

\item \textbf{Multi-Scale Integration:} The extension of graph methods to multi-scale HAR features demonstrates how network effects can be incorporated across different time horizons.

\item \textbf{Validation Framework:} We introduce rigorous validation methodologies to distinguish between genuine information content and parameter expansion effects:
\begin{itemize}
\item Random adjacency matrix validation for graph structure benefits
\item Pseudo-IV validation for implied volatility information content
\item Statistical testing frameworks for performance attribution
\end{itemize}

\item \textbf{Training Stability Analysis:} We identify and document critical training stability issues with QLIKE loss functions in neural network architectures, providing guidance for future implementations.
\end{enumerate}

\subsection{Practical Applications}

The findings have several practical applications:

\begin{itemize}
\item \textbf{Risk Management:} Improved volatility forecasts enable better Value-at-Risk and Expected Shortfall calculations.
\item \textbf{Portfolio Optimization:} More accurate volatility predictions support enhanced portfolio construction and rebalancing strategies.
\item \textbf{Derivatives Trading:} Better volatility forecasts can improve options pricing and volatility trading strategies.
\end{itemize}

\section{Limitations and Future Research}

\subsection{Current Limitations}

Several limitations should be acknowledged:

\begin{enumerate}
\item \textbf{Linear Models:} Our analysis focuses on linear model specifications, which may not fully exploit the potential of graph neural network architectures.

\item \textbf{Static Graphs:} The adjacency matrix is constructed using the full sample period, potentially missing time-varying network structures.

\item \textbf{Single Asset Universe:} The analysis is limited to Dow 30 stocks, and results may not generalize to other asset classes or market segments.
\end{enumerate}

\subsection{Future Research Directions}

Several avenues for future research emerge from this work:

\begin{enumerate}
\item \textbf{Dynamic Graph Construction:} Investigate time-varying graph structures that adapt to changing market conditions.

\item \textbf{Nonlinear Extensions:} Implement full Graph Neural Network architectures (GNNHAR) to capture nonlinear relationships and interactions.

\item \textbf{Alternative Graph Methods:} Explore other graph construction methods such as correlation networks, mutual information, or transfer entropy.

\item \textbf{Extended Asset Universe:} Apply the methodology to broader asset universes including international markets, bonds, and commodities.

\item \textbf{High-Frequency Extensions:} Extend the framework to intraday volatility forecasting using high-frequency data.
\end{enumerate}

\section{Potential Research Extensions}

Based on the current findings, several specific research questions warrant further investigation:

\subsection{Scale Effects in Graph Neural Networks}

\subsubsection{Research Question: Network Size and GNNHAR Performance}

An important question concerns whether Graph Neural Network HAR (GNNHAR) models demonstrate greater performance advantages over traditional GHAR models when applied to larger networks. Specifically:

\textbf{Hypothesis:} GNNHAR models may exhibit superior performance gains relative to GHAR models when the number of assets increases, due to their enhanced capacity to capture complex, nonlinear network interactions.

\textbf{Proposed Methodology:}
\begin{itemize}
\item Compare GNNHAR vs GHAR performance on Dow 30 (30 assets) vs S\&P 100 (100 assets)
\item Measure relative performance improvement: $\Delta = \frac{\text{MSE}_{\text{GHAR}} - \text{MSE}_{\text{GNNHAR}}}{\text{MSE}_{\text{GHAR}}}$
\item Test hypothesis: $\Delta_{\text{S\&P100}} > \Delta_{\text{Dow30}}$
\end{itemize}

\textbf{Economic Rationale:} Larger networks contain more complex interaction patterns that linear models may fail to capture, potentially favoring the nonlinear modeling capacity of GNNs.

\subsection{Training Stability and Loss Function Optimization}

\subsubsection{Research Question: QLIKE Loss Explosion Prevention}

The QLIKE loss function, while theoretically superior for volatility forecasting, can suffer from numerical instability during training:

$$\mathcal{L}_{\text{QLIKE}} = \frac{1}{NT}\sum_{t}\sum_{i}\left(\frac{v_{i,t}}{\hat{v}_{i,t}+\epsilon} - \log\frac{v_{i,t}}{\hat{v}_{i,t}+\epsilon} - 1\right)$$

\textbf{Observed Training Instability:} Our empirical results reveal severe QLIKE training instability in neural networks:
\begin{itemize}
\item GNNHAR-2L GLASSO QLIKE: Loss exploded to $2.02 \times 10^9$, Test MSE: 571.36
\item GNNHAR-3L GLASSO QLIKE: Loss reached $7.35 \times 10^5$, Test MSE: 24.32
\item GNNHAR+IV models with QLIKE: Similar explosion patterns observed
\end{itemize}

\textbf{Potential Solutions:}
\begin{enumerate}
\item \textbf{Adaptive Clipping:} Implement gradient clipping with adaptive thresholds based on loss magnitude
\item \textbf{Regularized QLIKE:} Add penalty terms to prevent extreme predictions:
$$\mathcal{L}_{\text{Reg-QLIKE}} = \mathcal{L}_{\text{QLIKE}} + \lambda \sum_i (\hat{v}_{i,t})^2$$
\item \textbf{Robust Epsilon Scheduling:} Use time-varying $\epsilon$ values during training to prevent division by zero
\item \textbf{Hybrid Loss Functions:} Combine MSE and QLIKE with adaptive weights:
$$\mathcal{L}_{\text{Hybrid}} = (1-\alpha_t) \mathcal{L}_{\text{MSE}} + \alpha_t \mathcal{L}_{\text{QLIKE}}$$
\item \textbf{Learning Rate Scheduling:} Implement aggressive learning rate decay when QLIKE loss exceeds threshold values
\item \textbf{Early Stopping:} Monitor loss explosion and revert to previous stable weights
\end{enumerate}

\subsection{Nonlinear Model Performance Evaluation}

\subsubsection{Research Question: GNNHAR vs GHAR Comparative Analysis}

While our current study focuses on linear GHAR models, the full GNNHAR implementation remains to be evaluated:

\textbf{GNNHAR Architecture:}
\begin{align}
Z_{:t-1} &= [V_{:t-1}, X_{:t-1}] \in \mathbb{R}^{N\times 6} \\
H^{(l+1)} &= \sigma(WH^{(l)}\Theta^{(l)}), \quad l = 0,\ldots,L-1 \\
\hat{v}_t &= H^{(L)}\theta^{(\text{out})}
\end{align}

\textbf{Key Research Questions:}
\begin{itemize}
\item Does the nonlinear GNN architecture provide significant improvements over linear GHAR?
\item What is the optimal network depth $L$ for volatility forecasting?
\item How does computational complexity scale with network size and depth?
\end{itemize}

\subsection{Implied Volatility Integration and Lag Structure}

\subsubsection{Research Question: IV-Based Lag Mitigation}

Traditional volatility models suffer from inherent lag due to their reliance on historical realized volatility. Implied volatility, being forward-looking, may help address this limitation:

\textbf{Lag Problem in Traditional Models:}
$$\hat{v}_t = f(v_{t-1}, v_{t-5:t-2}, v_{t-22:t-6})$$

\textbf{IV-Enhanced Lag Mitigation Hypothesis:}
$$\hat{v}_t = f(v_{t-1}, v_{t-5:t-2}, v_{t-22:t-6}, x_{t-1}, x_{t-5:t-2}, x_{t-22:t-6})$$

Our empirical findings provide preliminary evidence supporting this hypothesis. The consistent performance improvements from IV inclusion across all model specifications (6-7\% MSE reduction) suggest that forward-looking market information can indeed mitigate the inherent lag in volatility forecasting.

\textbf{Proposed Detailed Analysis:}
\begin{enumerate}
\item \textbf{Cross-Correlation Analysis:} Measure prediction lag using cross-correlation between actual and predicted volatility series
\item \textbf{Lead-Lag Relationship:} Compare lag structure between RV-only and RV+IV models using Granger causality tests
\item \textbf{Real-Time Performance:} Evaluate forecasting performance during high volatility periods when lag effects are most pronounced
\item \textbf{Impulse Response Analysis:} Study how quickly models respond to volatility shocks with and without IV information
\item \textbf{Rolling Window Analysis:} Examine temporal stability of lag reduction benefits across different market regimes
\end{enumerate}

\textbf{Economic Significance:} Reducing prediction lag has substantial practical implications for:
\begin{itemize}
\item Risk management systems requiring rapid volatility updates
\item High-frequency trading strategies sensitive to volatility timing
\item Options market making where volatility forecasting speed is critical
\item Portfolio rebalancing decisions during market stress periods
\end{itemize}

\subsubsection{Research Question: Optimal IV Integration Methods}

Several approaches exist for incorporating implied volatility into HAR-type models:

\textbf{Method 1: Additive Integration (Current Approach):}
$$E[v_t] = \alpha + V_{:t-1}\beta + X_{:t-1}\delta$$

\textbf{Method 2: Multiplicative Integration:}
$$E[v_t] = \alpha + V_{:t-1}\beta \odot (1 + X_{:t-1}\delta)$$

\textbf{Method 3: Attention-Based Integration:}
$$E[v_t] = \alpha + V_{:t-1}\beta + \text{Attention}(V_{:t-1}, X_{:t-1})\gamma$$

\textbf{Method 4: Regime-Dependent Integration:}
$$E[v_t] = \alpha + V_{:t-1}\beta + \mathbf{1}_{S_t} X_{:t-1}\delta$$
where $S_t$ represents market regime (low/high volatility, crisis/normal periods).

\subsection{Information Content Validation Framework}

\subsubsection{Research Question: Pseudo-IV Validation Test}

A fundamental question concerns whether the observed performance improvements from including implied volatility stem from genuine forward-looking information content or merely from parameter expansion. We propose a rigorous validation framework using pseudo-implied volatility.

\textbf{Pseudo-IV Generation Methodology:}

\textbf{Step 1: Statistical Characterization}
For each stock $i$ and each HAR component (daily, weekly, monthly), compute the empirical statistics of true implied volatility:
\begin{align}
\mu_{i,d} &= \mathbb{E}[x_{i,t-1}] \\
\sigma_{i,d} &= \text{Std}[x_{i,t-1}] \\
\mu_{i,w} &= \mathbb{E}[x_{i,t-5:t-2}] \\
\sigma_{i,w} &= \text{Std}[x_{i,t-5:t-2}] \\
\mu_{i,m} &= \mathbb{E}[x_{i,t-22:t-6}] \\
\sigma_{i,m} &= \text{Std}[x_{i,t-22:t-6}]
\end{align}

\textbf{Step 2: Pseudo-IV Generation}
Generate synthetic implied volatility series that preserve statistical properties but remove predictive information:
\begin{align}
\tilde{x}_{i,t-1} &\sim \mathcal{N}(\mu_{i,d}, \sigma_{i,d}^2) \\
\tilde{x}_{i,t-5:t-2} &\sim \mathcal{N}(\mu_{i,w}, \sigma_{i,w}^2) \\
\tilde{x}_{i,t-22:t-6} &\sim \mathcal{N}(\mu_{i,m}, \sigma_{i,m}^2)
\end{align}

\textbf{Step 3: Cross-Sectional Correlation Preservation}
To maintain realistic cross-asset relationships, apply Cholesky decomposition to preserve the correlation structure:
$$\tilde{X}_t = \mu + L \cdot Z_t$$
where $L$ is the Cholesky factor of the empirical IV correlation matrix and $Z_t \sim \mathcal{N}(0,I)$.

\textbf{Experimental Design:}

\textbf{Treatment Groups:}
\begin{enumerate}
\item \textbf{Control (RV-only):} HAR, GHAR, GNNHAR models using only realized volatility
\item \textbf{True IV:} HAR+IV, GHAR+IV, GNNHAR+IV models using actual implied volatility
\item \textbf{Pseudo IV:} HAR+PseudoIV, GHAR+PseudoIV, GNNHAR+PseudoIV using synthetic IV
\item \textbf{Random Features:} Models with additional random Gaussian features (same dimensionality as IV)
\end{enumerate}

\textbf{Hypothesis Testing:}
\begin{itemize}
\item \textbf{$H_1$ (Information Content):} $\text{MSE}_{\text{True IV}} < \text{MSE}_{\text{Pseudo IV}} \approx \text{MSE}_{\text{Random}}$
\item \textbf{$H_2$ (Parameter Effect):} $\text{MSE}_{\text{Pseudo IV}} < \text{MSE}_{\text{RV-only}}$
\end{itemize}

If $H_1$ holds, the performance gap between true and pseudo IV validates genuine information content. If only $H_2$ holds, improvements may be primarily due to increased model flexibility.

\textbf{Expected Results and Interpretation:}

\textbf{Scenario 1: Strong Information Content}
$$\text{MSE}_{\text{True IV}} \ll \text{MSE}_{\text{Pseudo IV}} \approx \text{MSE}_{\text{RV-only}}$$
This would indicate that IV improvements stem primarily from forward-looking information rather than parameter expansion.

\textbf{Scenario 2: Mixed Effects}
$$\text{MSE}_{\text{True IV}} < \text{MSE}_{\text{Pseudo IV}} < \text{MSE}_{\text{RV-only}}$$
This would suggest both information content and parameter flexibility contribute to performance improvements.

\textbf{Scenario 3: Primarily Parameter Effects}
$$\text{MSE}_{\text{True IV}} \approx \text{MSE}_{\text{Pseudo IV}} < \text{MSE}_{\text{RV-only}}$$
This would indicate that performance gains are mainly due to increased model complexity rather than genuine IV information.

\textbf{Statistical Testing Framework:}
\begin{enumerate}
\item \textbf{Diebold-Mariano Test:} Compare forecast accuracy between true IV and pseudo IV models
\item \textbf{Model Confidence Set:} Identify the set of statistically equivalent models
\item \textbf{Bootstrap Validation:} Generate confidence intervals for performance differences
\item \textbf{Cross-Validation:} Ensure robustness across different time periods
\end{enumerate}

\subsubsection{Extended Research Question: Neural Network Sensitivity to Pseudo-IV}

A critical extension of the pseudo-IV validation concerns whether neural networks exhibit different sensitivity to uninformative variables compared to linear models. This addresses the fundamental question of whether the increased complexity of GNNs makes them more susceptible to overfitting on noise.

\textbf{Hypothesis Framework:}

\textbf{$H_3$ (Neural Network Robustness):} 
$$\frac{\text{MSE}_{\text{GNNHAR+True IV}} - \text{MSE}_{\text{GNNHAR+Pseudo IV}}}{\text{MSE}_{\text{GNNHAR+True IV}}} > \frac{\text{MSE}_{\text{GHAR+True IV}} - \text{MSE}_{\text{GHAR+Pseudo IV}}}{\text{MSE}_{\text{GHAR+True IV}}}$$

If $H_3$ holds, neural networks demonstrate greater ability to distinguish between informative and uninformative variables.

\textbf{$H_4$ (Overfitting Susceptibility):}
$$\text{MSE}_{\text{GNNHAR+Pseudo IV}} < \text{MSE}_{\text{GNNHAR RV-only}}$$

If $H_4$ holds, it suggests neural networks may overfit to uninformative variables, potentially raising concerns about model robustness.

\textbf{Key Research Questions:}

\begin{enumerate}
\item \textbf{Comparative Sensitivity:} Do GNNs show larger performance gaps between true and pseudo IV than linear models?

\item \textbf{Overfitting Risk:} Can pseudo-IV improve GNN performance beyond RV-only baselines, indicating potential overfitting to noise?

\item \textbf{Architecture Dependence:} How does pseudo-IV sensitivity vary across different GNN depths (1L, 2L, 3L)?

\item \textbf{Training Stability:} Do GNNs with pseudo-IV exhibit different training dynamics compared to true IV?
\end{enumerate}

\textbf{Expected Outcomes and Implications:}

\textbf{Scenario A: Enhanced Discrimination ($H_3$ holds, $H_4$ fails)}
- GNNHAR shows larger true vs pseudo-IV gap than GHAR
- Pseudo-IV provides minimal or no improvement over RV-only GNN
- \textit{Implication:} Neural networks better distinguish signal from noise

\textbf{Scenario B: Overfitting Susceptibility ($H_4$ holds)}
- Pseudo-IV improves GNN performance beyond RV-only baseline
- May even approach true IV performance levels
- \textit{Implication:} GNNs prone to overfitting on uninformative variables

\textbf{Scenario C: Similar Sensitivity (Neither $H_3$ nor $H_4$ holds)}
- GNN and linear models show comparable pseudo-IV effects
- \textit{Implication:} Model complexity doesn't significantly affect noise sensitivity

\textbf{Methodological Significance:}

This extended validation addresses a crucial concern in deep learning applications to finance: whether increased model complexity leads to better signal extraction or merely sophisticated overfitting. The results will inform:

\begin{itemize}
\item \textbf{Model Selection:} When to prefer neural networks over linear models
\item \textbf{Regularization Strategies:} How to prevent overfitting in financial GNNs
\item \textbf{Feature Engineering:} Guidelines for including additional variables in neural volatility models
\item \textbf{Risk Management:} Understanding model reliability under different information conditions
\end{itemize}

\subsection{Empirical Research Program}

\textbf{Phase 1: Scale Analysis}
\begin{itemize}
\item Implement GNNHAR for both Dow 30 and S\&P 100
\item Compare relative performance improvements
\item Analyze computational scalability
\end{itemize}

\textbf{Phase 2: Information Content Validation}
\begin{itemize}
\item \textbf{Pseudo-IV Generation:} Create synthetic IV series preserving statistical properties but removing predictive content
\item \textbf{Linear Model Analysis:} Test HAR+IV, GHAR+IV models with true vs pseudo IV $\checkmark$ \textit{(Completed)}
\item \textbf{Neural Network Extension:} Test GNNHAR+IV models (1L, 2L, 3L) with pseudo-IV across different architectures
\item \textbf{Comparative Sensitivity Analysis:} Compare linear vs neural network responses to uninformative variables
\item \textbf{Statistical Testing:} Apply Diebold-Mariano tests and Model Confidence Sets
\item \textbf{Overfitting Assessment:} Evaluate whether pseudo-IV improves GNN performance beyond RV-only baselines
\end{itemize}

\textbf{Phase 3: Training Optimization}
\begin{itemize}
\item Develop robust QLIKE training procedures
\item Implement and test regularization techniques
\item Compare training stability across model specifications
\end{itemize}

\textbf{Phase 4: Nonlinear Extensions}
\begin{itemize}
\item Full GNNHAR implementation and evaluation
\item Architecture optimization (depth, width, activation functions)
\item Performance comparison with linear baselines
\end{itemize}

\textbf{Phase 5: IV Integration Optimization}
\begin{itemize}
\item Test alternative IV integration methods
\item Analyze lag reduction capabilities
\item Evaluate regime-dependent performance
\item Validate information content through pseudo-IV experiments
\end{itemize}

These research extensions would significantly enhance our understanding of graph-based volatility forecasting and provide practical insights for implementation in real-world trading and risk management systems.

\section{Future Research Program and Publication Strategy}

Based on the comprehensive findings from this study, we propose a structured research program that addresses fundamental questions in financial modeling methodology and practical applications. The research can be organized into three distinct but complementary papers, each making unique contributions to the literature.

\subsection{Paper 1: Methodological Innovation - Pseudo-Variable Validation Framework}

\subsubsection{Research Motivation and Discovery}

The genesis of this research stems from an unexpected finding in our simulation studies: randomly generated adjacency matrices occasionally outperformed GLASSO-constructed graphs, particularly in small networks. This counterintuitive result raised a fundamental question that pervades financial econometrics: \textbf{Are performance improvements in HAR model variants genuine information gains or merely artifacts of parameter expansion?}

\subsubsection{Core Methodological Contribution}

Many HAR model extensions in the literature claim superior forecasting performance by adding new variables (implied volatility, macroeconomic factors, sentiment indicators, etc.). However, these studies typically fail to control for the fundamental confound between information content and parameter count. Our pseudo-variable validation framework addresses this critical gap by:

\begin{enumerate}
\item \textbf{Pseudo-Graph Validation:} Using randomly generated adjacency matrices to isolate graph structure benefits from parameter expansion effects
\item \textbf{Pseudo-IV Validation:} Generating synthetic implied volatility that preserves statistical properties but removes predictive information
\item \textbf{Comparative Architecture Analysis:} Demonstrating differential sensitivity between linear and neural network models to uninformative variables
\end{enumerate}

\textbf{Key Empirical Findings:}
\begin{itemize}
\item Linear models: 82.4\% information content, 17.6\% parameter effects
\item Neural networks: 65.6\% information content, 34.4\% parameter effects
\item GNNs extract 63\% more genuine information but show 4× higher parameter sensitivity
\end{itemize}

\subsubsection{Broader Implications}

This methodology provides a rigorous framework for evaluating any proposed enhancement to financial forecasting models, addressing concerns about:
\begin{itemize}
\item Data mining and spurious correlations
\item Model complexity vs. genuine predictive power
\item Overfitting in high-dimensional financial models
\item Fair comparison across different model architectures
\end{itemize}

\subsection{Paper 2: Applied Research - Implied Volatility in Graph-Enhanced Models}

\subsubsection{Research Focus}

This paper examines the specific application of implied volatility as an exogenous variable in GNNHAR models, building on the validation framework established in Paper 1.

\textbf{Core Research Question:} How does the incorporation of forward-looking implied volatility information enhance volatility forecasting in graph neural network architectures?

\subsubsection{Key Contributions}

\begin{enumerate}
\item \textbf{Information Content Quantification:} Demonstrating that IV provides 7.63 percentage points of genuine information improvement in GNNHAR models
\item \textbf{Cross-Asset Information Propagation:} Analyzing how IV information spreads through graph structures
\item \textbf{Temporal Lag Mitigation:} Investigating whether forward-looking IV reduces the inherent lag in volatility forecasting
\item \textbf{Practical Implementation:} Providing guidance for incorporating options market data in operational risk management systems
\end{enumerate}

\textbf{Empirical Evidence:}
\begin{itemize}
\item GNNHAR+IV achieves 11.63\% total improvement over RV-only models
\item 65.6\% of improvement stems from genuine information content
\item Neural networks show superior information extraction compared to linear models
\end{itemize}

\subsection{Paper 3: Theoretical Investigation - Network Scale Effects in Financial GNNs}

\subsubsection{Research Motivation}

A puzzling observation emerged from comparing our Dow 30 results with existing literature: while Zhang Chao (2023) found GNNHAR superior to GHAR using S\&P 100 data, our Dow 30 analysis often shows GHAR outperforming GNNHAR. Combined with our earlier simulation studies using only six stocks, this suggests a critical scale dependency in graph neural network effectiveness.

\subsubsection{Core Hypothesis}

\textbf{Network Scale Hypothesis:} The relative performance advantage of GNNHAR over GHAR increases with network size, as larger graphs provide richer interaction patterns that neural networks can better exploit than linear models.

\textbf{Mathematical Framework:}
$$\text{Relative Advantage} = \frac{\text{MSE}_{\text{GHAR}} - \text{MSE}_{\text{GNNHAR}}}{\text{MSE}_{\text{HAR}}} \times 100\%$$

\textbf{Testable Predictions:}
\begin{itemize}
\item $\text{Advantage}_{\text{S\&P100}} > \text{Advantage}_{\text{Dow30}} > \text{Advantage}_{\text{Small-6}}$
\item Critical network size threshold where GNNs begin outperforming linear models
\item Computational efficiency trade-offs at different scales
\end{itemize}

\subsubsection{Research Design}

\begin{enumerate}
\item \textbf{Multi-Scale Analysis:} Compare GNNHAR vs GHAR performance across:
   \begin{itemize}
   \item Small networks (6-10 stocks)
   \item Medium networks (Dow 30)
   \item Large networks (S\&P 100, S\&P 500)
   \end{itemize}

\item \textbf{Performance Metrics:}
   \begin{itemize}
   \item Relative improvement percentages
   \item Information content extraction efficiency
   \item Computational cost per unit improvement
   \end{itemize}

\item \textbf{Theoretical Analysis:}
   \begin{itemize}
   \item Graph complexity measures (clustering, path length, density)
   \item Information propagation capacity
   \item Overfitting risk assessment
   \end{itemize}
\end{enumerate}

\subsubsection{Expected Contributions}

\begin{itemize}
\item \textbf{Practical Guidance:} When to use GNNs vs linear models based on network size
\item \textbf{Theoretical Understanding:} Why neural networks require minimum complexity to be effective
\item \textbf{Resource Allocation:} Cost-benefit analysis for implementing sophisticated models
\item \textbf{Scalability Insights:} How GNN advantages scale with market coverage
\end{itemize}

\subsection{Publication Timeline and Strategy}

\textbf{Phase 1 (Months 1-6):} Complete Paper 1 (Pseudo-Variable Framework)
\begin{itemize}
\item Target: Review of Financial Studies / Journal of Finance
\item Emphasis: Methodological innovation and broad applicability
\end{itemize}

\textbf{Phase 2 (Months 4-10):} Develop Paper 3 (Network Scale Effects)  
\begin{itemize}
\item Target: Journal of Financial Econometrics / Quantitative Finance
\item Emphasis: Theoretical insights and practical guidance
\end{itemize}

\textbf{Phase 3 (Months 8-12):} Finalize Paper 2 (IV Applications)
\begin{itemize}
\item Target: Journal of Banking \& Finance / Finance Research Letters
\item Emphasis: Applied research building on established methodology
\end{itemize}

\subsection{Broader Research Impact}

This research program addresses fundamental questions in financial modeling:

\begin{enumerate}
\item \textbf{Model Validation:} How to rigorously test whether model improvements are genuine
\item \textbf{Architecture Selection:} When complexity adds value vs. when simplicity suffices  
\item \textbf{Scale Dependencies:} How network effects vary with market size and connectivity
\item \textbf{Information Processing:} How different model types extract and utilize financial information
\end{enumerate}

The methodological innovations, particularly the pseudo-variable validation framework, have potential applications beyond volatility forecasting to any domain where distinguishing signal from noise is critical. This positions the research to make lasting contributions to both financial econometrics methodology and practical risk management applications.

\section{Conclusion}

This study provides compelling evidence for the benefits of incorporating graph structure and spillover effects in volatility forecasting models. The GHAR+IV model, which combines graph-enhanced HAR features with implied volatility information, achieves the best performance with a 10.2\% improvement over traditional HAR models.

Key conclusions include:

\begin{enumerate}
\item \textbf{Graph Structure Matters:} GLASSO-constructed adjacency matrices effectively capture cross-asset spillover effects, leading to consistent performance improvements. Random adjacency matrix validation confirms these benefits stem from genuine economic relationships rather than parameter expansion.

\item \textbf{Implied Volatility Contains Genuine Information with Model-Dependent Effects:} Our comprehensive pseudo-IV validation reveals that information content varies by model architecture: 82.4\% for linear models vs 65.6\% for neural networks. While GNNs extract more absolute information (7.63\% vs 4.67\%), they also show higher parameter sensitivity, confirming the information content hypothesis with important architectural nuances.

\item \textbf{Methodological Innovation and Architecture Insights:} The extended pseudo-IV validation framework provides a replicable methodology for evaluating both information content and model architecture effects. Key insight: neural networks show 2.96 percentage points advantage in information discrimination but 4× higher parameter sensitivity, informing optimal model selection strategies.

\item \textbf{Robust Performance:} The benefits of graph enhancement are robust across different loss functions and model specifications, with neural networks providing competitive but not superior performance to well-designed linear models.

\item \textbf{Training Stability Challenges:} QLIKE loss functions exhibit severe numerical instability in deeper neural networks, highlighting the need for robust optimization techniques in volatility forecasting applications.
\end{enumerate}

These findings contribute to the growing literature on network-based financial modeling and provide practical insights for volatility forecasting in interconnected financial markets. The methodology developed here offers a foundation for further research into graph-enhanced financial prediction models.

\section*{Acknowledgments}

The authors acknowledge the valuable contributions of the financial econometrics and network analysis communities in developing the theoretical foundations underlying this work.

\bibliographystyle{plainnat}
\bibliography{references}

\end{document}